\documentclass[11pt]{article}
\usepackage[margin=1in]{geometry}
\usepackage[T1]{fontenc}
\usepackage{lmodern,microtype,hyperref,enumitem,amsmath,amssymb,xcolor,booktabs}
\usepackage[numbers,sort&compress]{natbib}
\hypersetup{colorlinks=true,urlcolor=blue,citecolor=blue,linkcolor=blue}
\title{Q14: Discussion Questions}
\author{DLS 2026 --- Kukla}
\date{October 2026}
\begin{document}
\maketitle
\begin{enumerate}[leftmargin=1.5em]
\item If concepts are directions rather than neurons, what does that imply for anyone hoping to ``find the dog neuron''? How does superposition change the interpretability project?
\item Sparse autoencoders assume features are sparse. Where might that assumption fail, and what would the consequence be for the features they recover?
\item Factual knowledge can be localized to specific MLP weights and edited. What are the promises and the risks of model editing as a way to update or correct a deployed model?
\item The linear representation hypothesis is powerful but not universal. Propose a capability you would expect to be represented \emph{non}-linearly, and how you'd test it.
\item Given the weights-vs-context split, what is the right architecture for a model that should genuinely learn from a long-running user relationship? (Preview Q15.)
\end{enumerate}
\end{document}
